\section*{Bab \ref{ch:b1:stereochemistry-in-depth} --- Stereokimia Lanjut}

\begin{solution}{exo:b1:stereochemistry-in-depth:1}
\omterm{def:b1:stereochemistry-in-depth:isomers}{Isomer struktur}; \omterm{def:b1:stereochemistry-in-depth:enantiomers}{enantiomer}; \omterm{def:b1:stereochemistry-in-depth:enantiomers}{diastereomer} (\omterm{def:b1:stereochemistry-in-depth:isomers}{stereoisomer} yang tidak saling
merupakan bayangan cermin); dua \omterm{def:b1:stereochemistry-in-depth:isomers}{konformasi} dari satu senyawa;
\omterm{def:b1:stereochemistry-in-depth:enantiomers}{diastereomer} (\cip{R,S} adalah bentuk meso).
\end{solution}

\begin{solution}{exo:b1:stereochemistry-in-depth:2}
\ce{-OH} $>$ \ce{-NH2} $>$ \ce{-CH3} $>$ \ce{-H} (O, N, C, H).
\ce{-COOH} (O,O,O) $>$ \ce{-CHO} (O,O,H) $>$ \ce{-CH2OH} (O,H,H) $>$
\ce{-CH(CH3)2} (C,C,H). \ce{-Br} $>$ \ce{-Cl} $>$ \ce{-CCl3} (Cl,Cl,Cl)
$>$ \ce{-CH2Cl} (Cl,H,H): atom pertama menentukan sebelum tetangga-tetangganya.
\end{solution}

\begin{solution}{exo:b1:stereochemistry-in-depth:3}
2,3-Diklorobutana: dua pusat yang ekuivalen, \cip{R,R}, \cip{S,S}, dan
\cip{R,S} yang meso: tiga. 2-Bromo-3-klorobutana: dua pusat yang
berbeda, empat \omterm{def:b1:stereochemistry-in-depth:isomers}{stereoisomer}, tanpa bentuk meso. Pentana-2,4-diol: tiga,
dengan \cip{2R,4S} sebagai bentuk meso.
\end{solution}

\begin{solution}{exo:b1:stereochemistry-in-depth:4}
Pada C5: isopropenil (C,C,C, dengan ikatan rangkap dua dihitung dua
kali) $>$ C6 $>$ C4 $>$ H. C6 dan C4 sama-sama membawa (C,H,H); satu
langkah lebih jauh C6 menuju karbon karbonil (O,O,C) dan C4 menuju C3
(C,C,H): C6 menang. Dengan gugus isopropenil di depan dan H di belakang,
urutan isopropenil $\to$ C6 $\to$ C4 berputar searah jarum jam seperti
yang digambar: \cip{R}, karvon mint hijau; yang berbaji garis adalah
\cip{S}.
\end{solution}

\begin{solution}{exo:b1:stereochemistry-in-depth:5}
$\theta = 0$: kedua metil eklips, \qty{16.5}{kJ/mol}; $60^\circ$: gauche,
sekitar \qty{2.9}{kJ/mol} (minimum 2.8 pada $62^\circ$); $120^\circ$:
metil menutupi H, \qty{15.2}{kJ/mol}; $180^\circ$: anti, 0. Jika
energinya sama, ada dua \omterm{def:b1:stereochemistry-in-depth:conformations}{konformasi gauche} untuk satu anti: dua pertiga
gauche.
\end{solution}

\begin{solution}{exo:b1:stereochemistry-in-depth:6}
Rantai tegak dengan \ce{COOH} di atas dan \ce{CH3} di bawah; urutannya
$\ce{NH2} > \ce{COOH} > \ce{CH3} > \ce{H}$. Dengan \ce{NH2} di kiri dan
H di kanan, $\ce{NH2} \to \ce{COOH} \to \ce{CH3}$ berputar searah jarum
jam sebagaimana terlihat; H mendatar (ke arah pembaca), sehingga
deskriptornya dibalik: \cip{S}. Gugus amino berada di kiri.
\end{solution}

\begin{solution}{exo:b1:stereochemistry-in-depth:7}
\textit{cis}: pada karbon-karbon yang bersebelahan, satu ikatan ke atas
dan satu ke bawah di antara posisi aksial dan \omterm{def:b1:stereochemistry-in-depth:conformations}{ekuatorial} menghasilkan
aksial--ekuatorial pada kedua \omterm{def:b1:stereochemistry-in-depth:conformations}{kursi}. \textit{trans}: diekuatorial atau
diaksial. Pada \omterm{def:b1:stereochemistry-in-depth:conformations}{kursi} diekuatorial kedua metil saling gauche (satu
interaksi); isomer \textit{cis} mempunyai interaksi itu ditambah satu
metil \omterm{def:b1:stereochemistry-in-depth:conformations}{aksial}, sekitar \qty{5.7}{kJ/mol} lebih tinggi. Isomer
\textit{trans} lebih stabil.
\end{solution}

\begin{solution}{exo:b1:stereochemistry-in-depth:8}
$[\alpha] = 13.2/(2.00 \times 0.100) = +66.0^\circ$. Tabung \qty{1.00}{dm}:
$+6.60^\circ$. \omterm{def:b1:stereochemistry-in-depth:enantiomers}{Enantiomernya}: $-13.2^\circ$ dalam tabung \qty{2.00}{dm}.
\end{solution}

\begin{solution}{exo:b1:stereochemistry-in-depth:9}
\qty{25}{\celsius}: $\exp(5700/(8.314 \times 298.15)) = 10.0$, bagiannya
$10.0/11.0 = \qty{91}{\percent}$. \qty{-50}{\celsius}: $\exp(5700/(8.314
\times 223.15)) = 21.6$, \qty{96}{\percent}.
\end{solution}

\begin{solution}{exo:b1:stereochemistry-in-depth:10}
Pada etana keenam hidrogen sama: setiap \omterm{def:b1:stereochemistry-in-depth:conformations}{konformasi goyang} identik. Pada
butana metil belakang dapat berdiri berseberangan dengan metil depan
(anti) atau di sampingnya ($\pm 60^\circ$, gauche), tempat kedua metil
saling berdesakan. Populasi: anti 1, gauche $2\exp(-2.8/2.48) = 0.65$:
sekitar \qty{61}{\percent} anti dan \qty{39}{\percent} gauche.
\end{solution}

\begin{solution}{exo:b1:stereochemistry-in-depth:11}
C2 dan C4 stereogenik. C3 membawa H, OH, dan dua cabang yang identik
konstitusinya; keduanya hanya berbeda jika C2 dan C4 mempunyai
deskriptor yang berlawanan, dan hanya pada saat itu C3 stereogenik
(pseudoasimetris). Isomer-isomernya: \cip{2R,4R} dan \cip{2S,4S},
sepasang \omterm{def:b1:stereochemistry-in-depth:enantiomers}{enantiomer} (C3 tidak stereogenik); \cip{2R,4S} dengan kedua
susunan pada C3, dua bentuk meso. Seluruhnya empat.
\end{solution}

\begin{solution}{exo:b1:stereochemistry-in-depth:12}
$[\alpha]_{\text{sampel}} = 2.31/(1.00 \times 0.050) = +46.2^\circ$, dan
$2x - 1 = 46.2/66.0 = 0.700$: $x = 0.850$, \qty{85}{\percent} \omterm{def:b1:stereochemistry-in-depth:enantiomers}{enantiomer}
$(+)$ dan \qty{15}{\percent} \omterm{def:b1:stereochemistry-in-depth:enantiomers}{enantiomer} $(-)$.
\end{solution}

\begin{solution}{pb:b1:stereochemistry-in-depth:1}
\textbf{1.} Cincin sikloheksana: C1 membawa OH, C2 gugus isopropil, C5
gugus metil; C1, C2, dan C5 stereogenik.
\textbf{2.} $2^3 = 8$, empat pasang \omterm{def:b1:stereochemistry-in-depth:enantiomers}{enantiomer}.
\textbf{3.} Tidak ada susunan yang mempunyai bidang cermin: ketiga
substituen itu berbeda semuanya.
\textbf{4.} \ce{OH} $>$ C2 (C,C,H) $>$ C6 (C,H,H) $>$ H.
\textbf{5.} C1 (O,C,H) $>$ isopropil (C,C,H) $>$ C3 (C,H,H) $>$ H.
\textbf{6.} C6 dan C4 sama-sama (C,H,H); berikutnya, C6 menuju C1 (O,C,H),
C4 menuju C3 (C,H,H): C6 $>$ C4 $>$ \ce{CH3} $>$ H.
\textbf{7.} \cip{1S,2R,5S}.
\textbf{8.} Pada \omterm{def:b1:stereochemistry-in-depth:conformations}{kursi}, dua ikatan \omterm{def:b1:stereochemistry-in-depth:conformations}{ekuatorial} pada karbon-karbon yang
bersebelahan (C1, C2) mengarah satu ke atas dan satu ke bawah:
\textit{trans}; pada karbon 1 dan 3 cincin (C1 dan C5, melalui C6) kedua
ikatan \omterm{def:b1:stereochemistry-in-depth:conformations}{ekuatorial} mengarah ke sisi yang sama: \textit{cis}. Itulah
susunan mentol.
\textbf{9.} Ketiganya menjadi \omterm{def:b1:stereochemistry-in-depth:conformations}{aksial}.
\textbf{10.} Pada \omterm{def:b1:stereochemistry-in-depth:conformations}{kursi} yang dibalik, metil saja sudah memerlukan sekitar
\qty{5.7}{kJ/mol} dan gugus isopropil yang lebih besar lebih banyak lagi;
OH menambah biayanya sendiri yang lebih kecil, dan OH serta metil yang
\omterm{def:b1:stereochemistry-in-depth:conformations}{aksial}, pada karbon 1 dan 3 dan pada sisi yang sama, juga akan saling
berdesakan. Totalnya jauh di atas \qty{12}{kJ/mol}, rasio lebih dari
$\exp(12000/2479) \approx 130$ banding satu: mentol hampir seluruhnya
berada dalam \omterm{def:b1:stereochemistry-in-depth:conformations}{kursi} yang serba \omterm{def:b1:stereochemistry-in-depth:conformations}{ekuatorial}.
\textbf{11.} \omterm{def:b1:stereochemistry-in-depth:enantiomers}{Diastereomer}: hanya satu dari tiga pusat yang berbeda.
\textbf{12.} Dengan isopropil dan metil \omterm{def:b1:stereochemistry-in-depth:conformations}{ekuatorial}, OH-nya \omterm{def:b1:stereochemistry-in-depth:conformations}{aksial}.
\textbf{13.} \omterm{def:b1:stereochemistry-in-depth:enantiomers}{Diastereomer} berbeda dalam semua sifatnya (di sini \omterm{def:b1:intermolecular-forces-solvents:hbond}{ikatan
hidrogen} OH yang \omterm{def:b1:stereochemistry-in-depth:conformations}{aksial} atau \omterm{def:b1:stereochemistry-in-depth:conformations}{ekuatorial}); \omterm{def:b1:stereochemistry-in-depth:enantiomers}{enantiomer} hanya berbeda
terhadap pasangan-pasangan \omterm{def:b1:stereochemistry-in-depth:stereocentre}{kiral}.
\textbf{14.} $[\alpha] = -2.46/(1.00 \times 0.0500) = -49.2^\circ$, di
dalam rentang buku pegangan itu.
% ledger: alpha:l-menthol-range
\textbf{15.} $+49.2^\circ$; $0$.
\textbf{16.} $-1.97/0.0500 = -39.4^\circ$.
\textbf{17.} $\alpha = \ell c[\alpha]_{(-)}\bigl(x - (1 - x)\bigr) =
\ell c[\alpha]_{(-)}(2x - 1)$.
\textbf{18.} $x = \frac12(1 + \alpha/\alpha_{\text{acuan}})$ pada $\ell$
dan $c$ yang sama.
\textbf{19.} Ya: setiap senyawa lain yang aktif optis menambahkan
sukunya sendiri (\omterm{prop:b1:stereochemistry-in-depth:biot}{hukum Biot}). Pemeriksaan kemurnian kimia dengan metode
lain (kromatografi) diperlukan sebelum membaca rotasi sebagai komposisi.
\textbf{20.} $x = \frac12(1 + 1.97/2.46) = 0.900$.
\textbf{21.} $x = \frac12(1 + 2.40/2.46) = 0.988$.
\textbf{22.} Sampel kedua (\qty{98.8}{\percent}).
\textbf{23.} $-3.94^\circ$; kesalahan pembacaan yang sama lalu menjadi
bagian yang lebih kecil dari sudut itu.
\textbf{24.} \textbf{\qty{90.0}{\percent} $(-)$-mentol} (dan
\qty{10.0}{\percent} \omterm{def:b1:stereochemistry-in-depth:enantiomers}{enantiomernya}).
\end{solution}
